\documentclass{plantillaPPT}
\titulo{5}{Gradiente y Regla de la Cadena}
\begin{document}
\primeraDiapo

\begin{frame}{Derivada direccional máxima}
\framesubtitle{Vector gradiente}

Sabemos que si \(f:A\subseteq\mathbb{R}^n\to\mathbb{R}\) es diferenciable en \(\vec{a}\in A\) y \(\vec{v}\) es un vector unitario, entonces
\[
\frac{\partial f}{\partial \vec{v}}(\vec{a}) = \nabla f(\vec{a})\cdot\vec{v}
\]

\textbf{Pregunta: }¿Cuál es la dirección del vector \(\vec{v}\) que permite obtener una \textcolor{red}{derivada direccional máxima} en el punto \(\vec{a}\)?

Notar que
\[
\left|\frac{\partial f}{\partial \vec{v}}(\vec{a})\right|=|\nabla f(\vec{a})\cdot\vec{v}|\leq\|\nabla f(\vec{a})\|\,\text{(Por Cauchy-Schwarz)}
\]
Luego, \(\|\nabla f(\vec{a})\|\) es una cota superior de todas las derivadas direccionales de \(f\) en \(\vec{a}\).
    
\end{frame}

\begin{frame}{Derivada direccional máxima}
\framesubtitle{Vector Gradiente}

Si consideramos el vector \(\displaystyle\vec{v}=\frac{\nabla f(\vec{a})}{\|\nabla f(\vec{a})\|}\) (siempre que el vector gradiente no sea nulo) se tiene que:
\[
\frac{\partial f}{\partial \vec{v}}(\vec{a})=\nabla f\cdot\vec{v} = \nabla f(\vec{a})\cdot \frac{\nabla f(\vec{a})}{\|\nabla f(\vec{a}\|} = \|\nabla f(\vec{a})\|
\]

Por lo tanto, el vector \textcolor{red}{\(\nabla f(\vec{a})\)} proporciona \textcolor{red}{la dirección de mayor crecimiento de \(f\) en \(\vec{a}\)} y el valor máximo de crecimiento de la derivada direccional de \(f\) en \(\vec{a}\) es \textcolor{red}{\(\|\nabla f(\vec{a})\|\)}.\\[5pt]

\textcolor{blue}{Observación: } La \textcolor{blue}{mínima derivada direccional} de \(f\) en \(\vec{a}\) se alcanza en la dirección del vector \textcolor{blue}{\(-\nabla f(\vec{a})\)} y el valor mínimo de la derivada direccional de \(f\) en \(\vec{a}\) es \textcolor{blue}{\(-\|\nabla f(\vec{a})\|\)}.

\end{frame}

\begin{frame}{Problema 1}
\framesubtitle{Práctica 5}

1. La temperatura en cualquier punto de una plana rectangular ubicada en el plano \(xy\) es \(T(x,y)=3x^2+2xy\). La distancia se mide en metros.\\[5pt]
\((a)\) Calcule la máxima tasa de variación de temperatura en el punto \((3, -6)\) de la placa.\\[2pt]
\((b)\) Determine la dirección para la cual ocurre esta tasa de variación máxima en \((3,-6)\).
    
\end{frame}

\begin{frame}{Matriz Jacobiana}
\framesubtitle{Definición}

\begin{definicion}
Sea \(f:A\subseteq\mathbb{R}^n\to\mathbb{R}^m\) y \(\vec{a}=(a_1,a_2,...,a_n)\in A\) un punto interior. La matriz
    \[
    Jf(\vec{a})=
    \begin{pmatrix}
    \frac{\partial f_1}{\partial x_1}(\vec{a}) & \frac{\partial f_1}{\partial x_2}(\vec{a}) & \cdots & \frac{\partial f_1}{\partial x_n}(\vec{a}) \\
    \frac{\partial f_2}{\partial x_1}(\vec{a}) & \frac{\partial f_2}{\partial x_2}(\vec{a}) & \cdots & \frac{\partial f_2}{\partial x_n}(\vec{a}) \\
    \vdots & \vdots & \ddots & \vdots \\
    \frac{\partial f_m}{\partial x_1}(\vec{a}) & \frac{\partial f_m}{\partial x_2}(\vec{a}) & \cdots & \frac{\partial f_m}{\partial x_n}(\vec{a})
    \end{pmatrix}_{m\times n}
    \]
se llama \textcolor{red}{matriz Jacobiana} de \(f\) en \(\vec{a}\)
\end{definicion}
\end{frame}

\begin{frame}{La Diferencial}
\framesubtitle{Definición}

\begin{definicion}
    La transformación lineal
    \[
    \textcolor{red}{Df(\vec{a})}:\mathbb{R}^n\to\mathbb{R}^m
    \]
    \[
    \vec{x}=(x_1,x_2,...,x_n)\mapsto\textcolor{blue}{Df(\vec{a})(\vec{x})=Jf(\vec{a})_{m\times n}\cdot
    \begin{pmatrix}
        x_1\\x_2\\\vdots\\x_n
    \end{pmatrix}_{n\times 1}
    }
    \]
    se llama la \textcolor{red}{Diferencial} de \(f\) en \(\vec{a}\).
\end{definicion}
    
\end{frame}

\begin{frame}{Regla de la Cadena}
\vspace{-15pt}
\begin{teorema}
Sea \(g:A\subseteq\mathbb{R}^n\to\mathbb{R}^m\) y \(f:B\subseteq\mathbb{R}^m\to\mathbb{R}^p\). Si \(g\) es diferenciable en \(\vec{a}\) y \(f\) es diferenciable en \(g(\vec{a})\), entonces la función compuesta \(f\circ g\) es diferenciable en \(\vec{a}\) y se cumple la siguiente fórmula
\[
D(f\circ g)(\vec{a}) = Df(g(\vec{a}))\circ Dg(\vec{a})
\]

Del Álgebra Lineal, la composición de transformaciones lineales corresponde al producto de sus matrices asociadas. Entonces, la fórmula del teorema se traduce en
\[
\textcolor{red}{J(f\circ g)(\vec{a}) = Jf(g(\vec{a}))\cdot Jg(\vec{a})}.
\]

donde \(J(f\circ g)(\vec{a})\) es de tamaño \(p \times n\).
\end{teorema}
\end{frame}

\begin{frame}{Problema 2}
\framesubtitle{Práctica 5}

2. Sean \(f:\mathbb{R}^2\to\mathbb{R}^3\) y \(g:\mathbb{R}^3\to\mathbb{R}^2\) las funciones \(f(x,y) = (x+y, xy, x^2+y^2)\), \(g(x,y,z)=(x^2, y^2+z^3)\). \\[5pt]
Calcular \(J(f\circ g)(1,1,1)\).
    
\end{frame}

\begin{frame}{Problema 3}
\framesubtitle{Práctica 5}

3. Si \(f\) es una función diferenciable de \(x\) e \(y\), \(u=f(x,y)\), \(x=r\cos(\theta)\) y \(y=r\sen(\theta)\), muestre que\\
\((a)\)

\begin{align*}
    \frac{\partial u}{\partial x}=\frac{\partial u}{\partial r}\cos(\theta)-\frac{1}{r}\frac{\partial u}{\partial \theta}\sin(\theta)\\[10pt]
    \frac{\partial u}{\partial y}=\frac{\partial u}{\partial r}\sin(\theta)-\frac{1}{r}\frac{\partial u}{\partial \theta}\cos(\theta)
\end{align*}

\end{frame}

\begin{frame}{Problema 3}
\framesubtitle{Práctica 5}
\((b)\)
\[
\left(\frac{\partial u}{\partial x}\right)^2+\left(\frac{\partial u}{\partial y}\right)^2 = \left(\frac{\partial u}{\partial r}\right)^2+\frac{1}{r^2}\left(\frac{\partial u}{\partial \theta}\right)^2
\]
    
\end{frame}

\begin{frame}{Problema 4}
\framesubtitle{Práctica 5}

4. Sean \(f\) y  \(g\) dos funciones de una variable real de clase \(C^2\), y definamos
\[F(x,y) = f[x+g(y)].\]

Hallar las fórmulas correspondientes a todas las derivadas parciales de \(F\) de primero y segundo orden, expresadas en función de las derivadas de \(f\) y \(g\). Comprobar la relación
\[
\frac{\partial F}{\partial x}\frac{\partial^2 F}{\partial x\partial y}-\frac{\partial F}{\partial y}\frac{\partial^2 F}{\partial^2 x}=0
\]
    
\end{frame}

\end{document}